\section{Action and Contemplation in Dante’s Divine Comedy}

Titus Burckhardt, in his essay \textit{Because Dante Was Right}, argues that one of the main themes of the \emph{Commedia} is ``the reciprocal relationship between knowledge and will.''

\begin{quotex}
Knowledge of the eternal truths is potentially present in the human spirit or intellect, but its unfolding is directly conditioned by the will, negatively when the soul falls into sin, and positively when this fall is overcome. The different punishments in Purgatory that Dante describes can be regarded, not only as posthumous states, but also as stages in ascesis, that lead to the integral and primordial condition, in which knowledge and will -- or more precisely, knowledge of man's eternal goal and his striving after pleasure -- are no longer separated from one another.
\end{quotex}


Will is the motive force of the human being, and corresponds to the \emph{vita activa}. Knowledge refers to the mind and awareness, and corresponds to the \emph{vita contemplativa}. By using our will and directing our energies toward understanding and spiritual growth, we can bring about the ``unfolding'' of the eternal truths in ourselves. But conversely, it is our knowledge, or lack of knowledge, which determines how we apply our will. Thus, we may think we are on a spiritual path, but if we are being guided by faulty knowledge, we may only be chasing delusion.

All human beings desire happiness and pleasure, and exercise their will to varying degrees in pursuit of them. But in the `normal,' unenlightened state in which will and knowledge are separated, man lacks knowledge of what will truly bring him happiness and pleasure, and so he exercises his will in vain, seeking fulfillment and stability where there is none to be found.

The sinner and saint alike both strive for happiness and pleasure. But the sinner's pursuit is based on confusion, on ignorance of causality, and on mis-perception. Hence, he `misses the mark.' He seeks his happiness in physical pleasure; in entertainment; in wealth and fame and power; in information, theories, and other such pseudo-knowledge; and then he reaps disappointment when these worldly pursuits not only do not satisfy his longing, but intensify his suffering. Dante provides many graphic examples of this in the \emph{Inferno}.

The saint, or aspirant, sees the sinner's life as the error and folly that it is, for no true happiness is to be won that way. So rather than wander aimlessly on the horizontal plane of existence, and risk descending to the depths of despair and suffering, he wills the vector of his life upwards towards the transcendent and Absolute. There are multiple ways of doing this, as was explained in the article on the \emph{Fedeli d'Amore}. The relevant distinction here is between the ways of the \emph{vita activa}, such as knightly service in the Crusades, and the \emph{vita contemplativa}, such as was exemplified by Bernard of Clairvaux, who guides Dante through the highest heavenly spheres. Both ways are valid and complement one another. Bernard, after all, played a crucial role in the formation and structuring of the Templar Order.

The \emph{vita contemplativa} emphasizes the development of knowledge, and the \emph{vita activa} emphasizes the development of will. Regarding contemplation, Burckhardt tells us that ``the human spirit, by penetrating more and more deeply into the Divine Wisdom, becomes gradually transformed into it.'' Similarly, the human will, by submitting completely to the Divine Will, such as is exemplified in the total devotion of the Crusader Knight, gradually is transformed into it.

\begin{quotex} 
The will of him who knows God springs from the source of freedom itself. Thus real freedom of the will depends on its relationship with the truth, which forms the content of essential knowledge. Conversely, the highest vision of God, of which Dante speaks in his work, is in accord with the spontaneous fulfillment of the divine will. Here knowledge has become one with the divine truth and will has become one with the divine love; both qualities reveal themselves as aspects of Divine Being, the one static and the other dynamic.
\end{quotex}


Or as the Buddhist tradition teaches, the compassionate activity of an enlightened being is in no way contrived, but rather is the spontaneous and natural activity of the Buddha nature when it is no longer hindered by dualistic fixation and delusion. The primordial non-duality of emptiness and form is reflected in the non-duality of the clear, luminous mind of enlightenment, in which ``knowledge has become one with the divine truth,'' and the pure, compassionate form and activity of a Buddha, whose ``will has become one with the divine love.'' This is the state described by Dante in the final Canto of \emph{Paradiso}.

The path articulated so beautifully by Dante in the \emph{Commedia} is a path that leads from sin, error, and suffering, to repentance, conversion, ascesis, and redemption, and ultimately to wisdom, beatitude, and realization of the Absolute. He makes clear that there are many pitfalls along the way, and that even those who have attained to great heights can still miss the mark in subtle ways, such that the summit still eludes them. But, from the depths and heights of his great love, Dante produced for us a guide to the path, and a portrait of its ultimate goal.

For those interested in studying Dante in the context of Tradition, I recommend the following works:

\textsc{Henry Dwight Sedgwick}'s \textit{Dante} is an excellent introductory text.

\textsc{Titus Burckhardt} ``Because Dante is Right'' in \textit{Mirror of the Intellect}

\textsc{Christian Moevs}, \textit{The Metaphysics of Dante's Comedy}

I prefer Allen Mandelbaum's English translation of the Divine Comedy to the others I have seen.


\flrightit{Posted on 2010-06-30 by Will}

\begin{center}* * *\end{center}

\begin{footnotesize}
\begin{sffamily}

\texttt{Wayne Ferguson on 2017-10-27 at 09:10 said:}

Some of the recommended resources, available online:

\url{https://ia902709.us.archive.org/26/items/danteelementaryb00sedg/danteelementaryb00sedg.pdf}

\url{https://ia801607.us.archive.org/13/items/TBurckhardt/Titus\_Burckhardt-Mirror\_of\_the\_Intellect\_1987\_.pdf}

\url{https://english.duke.edu/sites/english.duke.edu/files/file-attachments/Moevs\_Ch4\_107-132\_0.pdf}

\hfill

\end{sffamily}
\end{footnotesize}

